\documentclass[11pt]{article}
\usepackage[T1]{fontenc}
\usepackage[margin=2.4cm]{geometry}
\usepackage{amsmath,amssymb}
\usepackage{braket}
\usepackage{enumitem}
\usepackage[colorlinks=true,linkcolor=blue,urlcolor=blue]{hyperref}
\setlength{\parindent}{0pt}
\setlength{\parskip}{0.6\baselineskip}
\newenvironment{referee}{\begin{quote}\itshape}{\end{quote}}

\begin{document}

\begin{flushright}
Vienna, \today
\end{flushright}

Dr.~Bill Munro\\
Associate Editor, Physical Review A

\textbf{Re: XU11228A}, ``Answer partitions and oracle access determine quantum
query complexity'' by Karl Svozil\\
(revised title: ``Answer Partitions, Oracle Response Spectra, and Exact Quantum Query Advantage'')

Dear Dr.~Munro,

I thank both referees for careful reports. Both ask for a clearer separation
between known and new results, and the second referee raises the numerical
status of part of the four-bit census. I have revised the manuscript
accordingly. The main changes are:

\begin{enumerate}[leftmargin=*]
\item \textbf{The four-bit census is now exact.} For the 43 orbit
representatives (6040 partitions) that were previously supported only by
numerical SDP infeasibility, the revised manuscript gives rational dual
certificates of the two-query Barnum--Saks--Szegedy program. Each certificate
bounds the optimal two-query worst-case success by a value between 0.9001 and
0.9862, which proves \(Q_E=3\). The certificates are lists of integers
verified in exact arithmetic by a standalone verifier included in the
supplement; the validity argument is a three-line weak-duality chain,
Eq.~(17). All 6435 balanced four-bit partitions are therefore classified
exactly (Table~III).
\item \textbf{Prior work versus new results.} A new Table~I lists, topic by
topic, what was known and what is added. The introduction now names the main
new result (the response-spectrum criterion, Theorem~1 and Proposition~4) and
the second (the exact four-bit classification), and states that
Propositions~1 and~2 are standard.
\item \textbf{Three-bit values credited.} The abstract, the caption of
Table~II, and Sec.~IV now state that the exact three-bit complexities and the
majority construction are those of Montanaro, Jozsa, and Mitchison.
\item \textbf{Bounded error.} New Proposition~5 treats odd cyclic responses
with bounded error, includes the second referee's construction, and proves
that the improved one-query success \((3+\cos\frac{\pi}{d})/4\) is optimal.
\item \textbf{Access conventions.} A new paragraph before Theorem~1 states the
access assumptions and the phase convention.
\item \textbf{Generalization and structure.} Proposition~4 extends Theorem~1
to a complete classification of one-query-exact tasks for point oracles on any
number of bits and reduces, for the XOR oracle, to the characterization of
Chen, Ye, and Li [Phys.\ Rev.\ A \textbf{101}, 022325 (2020)]. The cost of
identification with an inaccessible witness is now determined exactly
(Proposition~6). The manuscript is now organized as a regular article with sections
and numbered statements, and the title names the main result. To make the
paper easier to follow, the strategy common to all results is now stated in
four steps at the beginning of the Discussion (Sec.~VI) and summarized in
prose in the Introduction.
\end{enumerate}

Point-by-point replies follow. Section, equation, table, and statement numbers
refer to the revised manuscript. A list of changes is given at the end, and a
marked-up PDF is attached. References to ``Ref.~[$n$] of the submitted
version'' use the numbering of the submitted manuscript, as in the reports.

\section*{Reply to the First Referee}

\begin{referee}
I am not very clear about the contribution relative to prior work. The
manuscript should state more explicitly what is already known and what is
established for the first time here, particularly in relation to Refs.~[13]
and [14]. [\dots] A concise comparison paragraph or a small table
distinguishing prior results from the present additions would be sufficient.
\end{referee}

I agree, and I have added both. The introduction now ends with a paragraph
naming the new results, and Table~I separates prior results from the present
additions. With respect to the two references named:

\begin{itemize}[leftmargin=*]
\item Montanaro, Jozsa, and Mitchison (Ref.~[13] of the submitted version)
give explicit optimal exact algorithms for every function on up to three
bits, and numerical SDP values for all functions on four bits, which they
note do not by themselves certify exactness. The revised manuscript states
that the three-bit values in Table~II are theirs and adds no new three-bit
values. For four bits it adds exact rational dual certificates for every
balanced partition whose two-query infeasibility was previously only
numerical (Sec.~IV, Table~III).
\item Ye, Yang, and Huang (Ref.~[14] of the submitted version) study exact
two-query algorithms using the four-bit numerical classification. The present
certificates make the balanced part of that classification rigorous.
\end{itemize}

The results that are new, to my knowledge, are the response-spectrum
criterion for an arbitrary response unitary (Theorem~1, Proposition~4), its
consequences for reversible completions and phase dressing (Corollary~1), the
bounded-error statement for odd cyclic responses (Proposition~5), the exact
four-bit certificates, and the exact costs of the witness channel
(Proposition~6). Propositions~1 and~2 are standard discrimination conditions
and convex duality, stated in the form used later.

\section*{Reply to the Second Referee}

\begin{referee}
(1) [\dots] Ref.~[13] [\dots] also provides explicit optimal exact algorithms
for all Boolean functions on up to three bits. Its Sec.~6.2 includes the same
majority construction described here. Thus, the exact three-bit query
complexities are already known. The author should acknowledge this directly
when presenting Table~I and in the abstract.
\end{referee}

The referee is right. The abstract now says that the paper reorganizes the
known exact complexities of the 35 balanced three-bit partitions. The caption
of the three-bit table (now Table~II) states that its values are given in
Ref.~[13], and the text in Sec.~IV says that the contribution for three bits
is only their organization as balanced partitions and that the majority
algorithm is the one in Sec.~6.2 of Ref.~[13].

\begin{referee}
(2) More generally, the manuscript should identify its main original
scientific result earlier. The dependence of query complexity on the oracle
interface is already discussed in Ref.~[1]. Proposition~1 is also
acknowledged to follow from standard discrimination conditions. The
response-spectrum theorem appears to be the most concrete contribution. Its
significance would be easier to assess if the introduction clearly
distinguished this result from the background examples.
\end{referee}

I agree. The first paragraph of the introduction now states that the
dependence on the oracle model is known in general, citing Ref.~[1] (W\'ojcik
and Chhajlany) together with Kashefi \emph{et al.}, Calude, and Regev and
Schiff. The second paragraph names the response-spectrum criterion as the main
new result, identifies the exact four-bit classification as the second, and
describes Propositions~1 and~2 as framework. Table~I makes the same separation
item by item.

The response-spectrum result has also been strengthened. Proposition~4 shows
that for point access through any response unitary \(V\) on any number of
bits, one query resolves a dictator exactly if and only if \(0\) lies in the
convex hull of the spectrum of \(V\), and a two-bit parity if and only if
\(-1\) is an eigenvalue of \(V\); no other task is one-query exact. For the
XOR oracle this is the theorem of Chen, Ye, and Li; for an odd cyclic response
the one-query class reduces to the dictators; and for \(V=-I\), the sign oracle,
it contains the two-bit parities but no dictator. The proof of Theorem~1 now uses
an explicit dual witness, \(Z=2I_2\otimes(I+V)(I+V)^\dagger\) [Eq.~(11)],
which is positive definite exactly when \(-1\notin\operatorname{spec}V\).

\begin{referee}
(3) The even-odd distinction in Corollary~1 concerns exact computation. A
short discussion of bounded-error computation would help clarify its
significance. [\dots] Hence one query still gives a bounded-error quantum
advantage [\dots]. Please discuss this distinction.
\end{referee}

I thank the referee for this point; it is now part of the paper. Section~III
reproduces the referee's construction (success 1 on constant and \(3/4\) on
balanced instances, against \(1/2\) for one classical query) and then proves a
sharper statement, Proposition~5. For odd \(d\), a probe spread equally over
the two eigenvectors of \(X_d\) whose eigenvalues are adjacent to \(-1\)
solves Deutsch's problem with one query and zero error. The answer is
conclusive with probability \((1+\cos\frac{\pi}{d})/2\) on every instance, and
guessing on the inconclusive outcome gives worst-case success
\((3+\cos\frac{\pi}{d})/4\). For \(d=3\) these values are \(3/4\) and \(7/8\).
The value is optimal: for any probe, the cell-averaged output states differ
only in an off-diagonal block bounded by \(\lVert I-V\rVert^2\), which limits
every one-query protocol to success \(\frac12+\frac18\lVert I-V\rVert^2\), equal
to \((3+\cos\frac{\pi}{d})/4\) for odd \(d\). The same bound reaches 1 only when
\(-1\) is an eigenvalue of \(V\), which gives a second proof of Theorem~1. One
classical query never yields a conclusive zero-error answer. The text now says explicitly that the odd cyclic response
removes the exact advantage but not the bounded-error one, and the abstract
mentions the \(7/8\) value.

\begin{referee}
(4) The access assumptions underlying Theorem~1 should be stated explicitly.
[\dots] the sufficiency argument assumes that an appropriate eigenstate of the
chosen \(V\) can be prepared at no query cost. Presumably, \(V\) is known, and
arbitrary, instance-independent operations are allowed. There is also a phase
convention worth explaining. Replacing \(V\) by \(e^{i\theta}V\) can change the
condition in the theorem. In Eq.~(9), this phase becomes relative between
address sectors and can therefore be observable. [\dots]
\end{referee}

The referee's reading is correct, and the assumptions are now stated in a
paragraph ``Access conventions'' directly after the definition of the
controlled-response family [now Eq.~(10)]. The family is specified and \(V\) is
known; any instance-independent operation, including the preparation of
eigenvectors of \(V\), is free; the oracle may not be controlled or inverted
unless that access is specified separately. The paragraph explains that
\(e^{i\theta}V\) multiplies the oracle by the instance-dependent phase
\(e^{i\theta b_x}\), which is relative between address sectors and observable,
and that this distinguishes the specified oracle from access to an unspecified
unitary or channel known only up to a global phase. Corollary~1 now uses the
phase convention as an example: \(X_3\) and \(e^{i\pi/3}X_3\) act identically
on every computational-basis input up to a phase, yet the second has
eigenvalue \(-1\) and gives \(Q_E=1\). The corollary also compares two
reversible permutation completions of the same blank-register table, the
transposition \(\ket0\leftrightarrow\ket1\) (spectrum \(\{1,1,-1\}\),
\(Q_E=1\)) and the cyclic shift (\(Q_E=2\)).

\begin{referee}
(5) The numerical status should be reflected consistently in the
supplementary data. For the 43 unresolved four-bit orbits, the CSV column
``Q\_lower\_bound'' and the JSON field ``lower\_bound'' contain the value 3.
However, the rigorous bounds supplied for these cases are \(2\le Q_E\le3\).
[\dots] Please separate the proved lower bound from the numerical candidate.
Exact certificates would strengthen the census, but repeating SCS at a tighter
tolerance does not establish these missing lower bounds, as the manuscript
correctly notes.
\end{referee}

The referee was right that the submitted supplement mixed proved and
numerical bounds. Following the suggestion, the missing lower bounds are now
proved by exact certificates rather than by further numerics. I used the
bounded-error form of the two-query SDP, whose dual is well posed, rather than
an exact-feasibility test. To a numerical dual solution I added
\(3\varepsilon\), \(2\varepsilon\), and \(\varepsilon\) times the identity to
\(Y_1\), \(Y_2\), and \(Y_3\) (\(\varepsilon=2^{-20}\)), which raises every
constraint matrix by \(\varepsilon I\), and rounded all entries to dyadic
rationals with denominator \(2^{40}\); all twelve matrix inequalities and the
conditions on \(\lambda\) were then verified in exact integer arithmetic. All 43 certificates are valid, so \(Q_E=3\) for
these orbits is now a theorem (Sec.~IV, Table~III). Their numerical optima,
0.9000--0.9861, agree with the values reported by Montanaro, Jozsa, and
Mitchison, and the 15 orbits settled by trees give dual value 1, as they
must.

The supplement is now submitted as APS Supplemental Material. In its principal
files (\nolinkurl{data/results-m4-orbits.csv}, \nolinkurl{witnesses-m4/}, and the
cross-check \nolinkurl{data/m4-mjm-crosscheck.csv}) the 43 rows formerly labelled
\nolinkurl{numerical_candidate} now carry \nolinkurl{exact_dual_certificate}, and
every row states the source of its lower bound in a new field
(\nolinkurl{Q_lower_bound_source}: trivial, the exact one-query obstruction,
or a rational dual certificate), together with the certified two-query success
bound and the location of the certificate. A new folder
\nolinkurl{dual-certificates/} contains the 43 certificates, keyed by the
partition numbers of the CSV file, and a verifier that uses only the Python
standard library. The raw SCS rerun files are kept unchanged as the numerical
record; a note in that folder states that their lower bounds were numerical
when written and are superseded by the certificates.

\begin{referee}
(6) A few presentation details should be corrected. The supplementary PDF and
README use a different manuscript title. The notation \(D\) should be
explicitly related to the previously introduced classical quantity \(R_0\).
The term ``kinematic partition'' is introduced but does not seem necessary for
the subsequent analysis.
\end{referee}

All three are corrected. The supplementary PDF, the README, and the Lean
documentation now carry the revised manuscript title and mention the title
under which it was submitted. Section~II now states that \(D=R_0\) in the worst-case
convention used in the paper, because fixing the random string of a
zero-error protocol with a worst-case query bound leaves a deterministic
protocol with the same bound, and notes that zero-error cost measured in
expectation can be smaller. The term ``kinematic partition'' has been
removed.

\section*{List of changes}

\begin{itemize}[leftmargin=*]
\item \textbf{Title} changed to ``Answer Partitions, Oracle Response Spectra, and Exact Quantum Query Advantage''.
\item \textbf{Abstract} rewritten: main result first; three-bit values
described as known; four-bit classification exact; bounded-error value
\(7/8\) added.
\item \textbf{Format}: \texttt{pra} class, sections, theorem environments with
proofs.
\item \textbf{Sec.~I (Introduction)}: prior work on oracle-model dependence
cited (W\'ojcik--Chhajlany, Kashefi \emph{et al.}, Calude, Regev--Schiff);
relation to earlier state-partition work stated; new paragraph summarizing,
in plain prose, the four-step strategy used throughout the paper, to make the
argument easier to follow; paragraph naming the new results; new Table~I
(prior work versus this work); closing paragraph on quantum parallelism
shortened, since the new summary covers its first part.
\item \textbf{Sec.~II}: term ``kinematic partition'' removed; \(D\) defined
and related to \(R_0\); Proposition~1 with proof; new Proposition~2 (dual
certificate) and its relation to adversary progress measures; remark on
commuting oracle families.
\item \textbf{Sec.~III}: Deutsch citation corrected to include the exact
one-query version (Cleve \emph{et al.}); new Proposition~3 (sign oracle
cannot resolve partitions that separate a truth table from its complement);
access-conventions paragraph; Theorem~1 proved with an explicit dual witness;
Corollary~1 rewritten (transposition versus cyclic shift, phase-dressed
shift, general \(X_d\) with \(d\ge2\)); new Proposition~4 (complete one-query
classification for every response unitary); new Proposition~5 (optimal
bounded-error solution for odd cyclic responses, \(d\ge3\)), with the
referee's construction; parity paragraph with explicit dual witness, degree
bound for every response, and an open question.
\item \textbf{Sec.~IV}: three-bit values credited to Montanaro--Jozsa--Mitchison
(caption of Table~II and text); truth-table bits in Table~II written \(b_i\),
as in the text; one-query lower bounds from Proposition~4;
two-query lower bounds from rational dual certificates [Eqs.~(16)--(17)];
new Table~III for four bits; all counts now exact (4847 partitions with an
exact advantage, 1588 without); real-degree remark; new paragraph describing
the Lean~4 formalization in the supplement and its trust assumptions.
\item \textbf{Sec.~V}: witness channel written as dephasing followed by the
oracle; register of the reduced state stated; new Proposition~6 (exact
identification cost \(\min\{k+1,n\}\) with a witness of \(k\) address bits).
\item \textbf{Sec.~VI (Discussion)}: new opening that states the strategy of
the paper in four steps (answer partition, oracle access, probe and one call,
orthogonality across cells), each with a short intuitive explanation and a
pointer to where it is worked out, followed by the Deutsch example and a
remark on several calls and bounded error; added to make the paper easier to
follow. The paragraph after it now says which results vary which step. New
remark that the Deutsch and
Bernstein--Vazirani separations use only Clifford operations, so stabilizer
simulability and query advantage are distinct (with Gottesman,
Aaronson--Gottesman, and Veitch \emph{et al.}); Shor and Grover cited;
asymptotic remark sharpened.
\item \textbf{References}: added Chen--Ye--Li, Ye--Li, Kashefi \emph{et al.},
Calude, Regev--Schiff, Ambainis (adversary), Cleve \emph{et al.},
Deutsch--Jozsa, Naseri \emph{et al.}, Gottesman, Aaronson--Gottesman,
Veitch \emph{et al.}, Shor, Grover, and earlier work on state partitions;
Aaronson's open-problems survey removed.
\item \textbf{Use of AI tools} disclosed according to the APS policy: tools and
versions, the tasks they assisted with, and how the results were verified, in
the Acknowledgments and, for the computations, in the Methods paragraph of
Sec.~IV.
\item \textbf{Data availability}: the supplement is now provided as APS
Supplemental Material, with the standard citation in the manuscript.
\item \textbf{Supplemental Material}: new folder with the rational dual
certificates and their exact verifier; lower-bound fields labeled by source,
and numerical-candidate labels replaced, in the principal files; new section
on the certificates in the supplementary PDF; title corrected; the term
``kinematic'' removed from the Lean sources.
\end{itemize}

Sincerely,

Karl Svozil

\end{document}
